% 张三慧V1第9章《温度和气体动理论》独立候选解。
% 主教材PDF物理页335--337；原图9.18临界点精度仍待核验。
\section*{第9章\quad 温度和气体动理论：独立解析}
\paragraph{习题 9.1}
定容气体温度计理想模型给出 \(p/T=\mathrm{const}\)，
三相点温度 \(T_*=273.16\,\mathrm K\)。
故
\[
 \boxed{p(373.15)=6650\frac{373.15}{273.16}
 \approx9.08\times10^3\,\mathrm{Pa}},
\quad
 \boxed{t_C=273.16\frac{2200}{6650}-273.15
 \approx-182.8^\circ\mathrm C}.
\]
低温真实气体偏离理想模型时必须外推压强至零密度。

\paragraph{习题 9.2}
设 \(u=T-T_s\)，则
\(-du/dt=Au\)，变量分离
\(du/u=-A\,dt\)，积分
\(u(t)=(T_1-T_s)e^{-At}\)，所以
\[
 \boxed{T(t)=T_s+(T_1-T_s)e^{-At}}.
\]
冬日数据 \(T_s=-10,T_1=25,T(40)=20\)
（温差单位均为摄氏度）满足
\(e^{-40A}=30/35=6/7\)，
\(A=(1/40)\ln(7/6)\,\mathrm{min^{-1}}\)。
下降至 \(15^\circ\mathrm C\) 时
\(e^{-At}=25/35=5/7\)，
\[
 \boxed{t=40\frac{\ln(7/5)}{\ln(7/6)}
 \approx87.3\,\mathrm{min}}.
\]
即从停暖算起，而不是在前40min之后再过87.3min。

\paragraph{习题 9.3}
将轮胎视为圆环管，胎圈中心线半径
\(R=(.7112-.030)/2=.3406\,\mathrm m\)，
内胎截面半径 \(r=.015\,\mathrm m\)。
胎内有效体积
\(V_t=2\pi^2Rr^2\approx1.513\,\mathrm L\)。
每次打气筒输送的环境压强、环境温度气体体积
\(V_p=\pi r_p^2 L_p
 =\pi(.015)^2(.30)\approx.2121\,\mathrm L\)。
假定轮胎初态 \(p_0,T_0=270.15\,\mathrm K\)，
泵每次从同一环境吸入理想气体且无回流，
最后 \(T_f=280.15\,\mathrm K\)，
粒子数守恒：
\[
 \boxed{p_f=p_0\frac{T_f}{T_0}
 \left(1+\frac{20V_p}{V_t}\right)
 \approx3.94\,p_0
 \approx4.00\times10^5\,\mathrm{Pa}}.
\]
这是\emph{绝对压强}；表压约 \(2.94p_0\)。
车胎内圈实际体积和活塞死腔会改变数值，
需对原书图/建模口径作二次审核。

\paragraph{习题 9.4}
由理想气体
\(p=\rho RT/M=nk_BT\) 得
\[
 \boxed{T=\frac{pM}{\rho R}
 =\frac{(.18)(.029)}
 {(3.2\times10^{-6})(8.314)}
 \approx196\,\mathrm K},
\]
\[
 \boxed{n=\frac p{k_BT}
 \approx6.65\times10^{19}\,\mathrm{m^{-3}}}.
\]
密度取的是气体质量密度而不是分子数密度。

\paragraph{习题 9.5}
球内外压强相同，热空气在内
密度 \(\rho_i=\rho_oT_o/T_i\)，
外界密度
\(\rho_o=pM/(RT_o)\approx1.205\,\mathrm{kg/m^3}\)。
临界升力
\((\rho_o-\rho_i)gV=mg\)，
所以
\[
 \boxed{T_i=\frac{T_o}{1-m/(\rho_o V)}
 \approx356.5\,\mathrm K\approx83.4^\circ\mathrm C}.
\]
这里 \(m=4500\,\mathrm{kg}\) 为\emph{气球壳和负载}
而不是再加热空气自身的重量；
热空气的重量已由浮力密度差计算。
假设球壳体积不变、与外界等压。

\paragraph{习题 9.6}
分子数满足
\(N=pV/(k_BT)\)，
\(p=1.00\times10^{-18}\times101325\,\mathrm{Pa}\)，
\(V=10^{-6}\,\mathrm{m^3}\)，
\(T=293.15\,\mathrm K\)，
\[
 \boxed{\langle N\rangle\approx25.0}.
\]
这只是在热平衡统计模型中的\emph{平均分子数}，
单次观察是整数且有涨落。

\paragraph{习题 9.7}
\[
 \boxed{p=nk_BT
 =(10^{10})(1.380649\times10^{-23})(10^4)
 \approx1.38\times10^{-9}\,\mathrm{Pa}}.
\]
未计星际氢的离子化、辐射及磁压。

\paragraph{习题 9.8}
平衡微元气柱满足 \(dp/dy=-\rho g\)。
用 \(\rho=Mp/(RT)\)、\(T(y)=T_0-\alpha y\)，得
\[
 \frac{dp}{p}=-\frac{Mg}{R(T_0-\alpha y)}\,dy.
\]
由地面 \(p_0\) 积分：
\[
 \boxed{p(y)=p_0
 \left(1-\frac{\alpha y}{T_0}\right)^{Mg/(\alpha R)}},
 \qquad 0<\alpha y<T_0.
\]
当 \(\alpha\to0\)，
\(\lim_{\alpha\to0}\alpha^{-1}\ln(1-\alpha y/T_0)
 =-y/T_0\)，故
\(\boxed{p=p_0e^{-Mgy/(RT_0)}}\)。
以峰高 \(h=8848\,\mathrm m\) 作约定近似，
\(\alpha=.006\,\mathrm{K/m}\)：
\[
 \boxed{T(h)=273-.006(8848)
 \approx219.9\,\mathrm K=-53.2^\circ\mathrm C},
\]
\[
 \boxed{p(h)\approx2.96\times10^4\,\mathrm{Pa}
 \approx0.292\,\mathrm{atm}}.
\]
这是源书的恒定温度直减率模型，
不作为真实每日珠峰气象预报。

\paragraph{习题 9.9}
令 \(\boldsymbol w=\boldsymbol v_1-\boldsymbol v_2\)。
两分子独立、各分量服从方差
\(\sigma^2=k_BT/m\) 的正态分布，
则 \(w_x,w_y,w_z\) 各自的方差为
\(2\sigma^2\)，
故相对速率 \(|\boldsymbol w|\)
服从温度为 \(2T\) 所对应参数的麦克斯韦分布。
单分子的平均速率
\(\bar v=\sqrt{8k_BT/(\pi m)}\)，
因此
\[
 \boxed{\bar u=\sqrt{\frac{8k_B(2T)}{\pi m}}
 =\sqrt2\,\bar v}.
\]
注意本式是\emph{相对速率的平均值}，
不是均方根之差。

\paragraph{习题 9.10}
令 \(f(v)\) 为任意归一化非负速率分布，
定义 \(\bar v=\int vf(v)\,dv\)。
由于
\[
 0\le\int_0^\infty(v-\bar v)^2f(v)\,dv
 =\overline{v^2}-(\bar v)^2,
\]
立即得
\[
 \boxed{\sqrt{\overline{v^2}}\ge\bar v}.
\]
只有速率不涨落、全部粒子速率相同时才取等号。

\paragraph{习题 9.11}
硬球气体的平均自由程
\[
 \lambda=\frac1{\sqrt2\pi d^2n},\qquad
 n=p/(k_BT).
\]
取标准态 \(T=273.15\,\mathrm K,p=101325\,\mathrm{Pa}\)，
\(d=3.8\times10^{-10}\,\mathrm m\)，
\[
 \boxed{\lambda\approx5.80\times10^{-8}\,\mathrm m}.
\]
氮气平均速率
\(\bar v=\sqrt{8RT/(\pi M)}
 \approx4.54\times10^2\,\mathrm{m/s}\)，
故平均碰撞间隔
\[
 \boxed{\tau\simeq\lambda/\bar v
 \approx1.28\times10^{-10}\,\mathrm s}.
\]
若采用碰撞频率的不同平均定义，
须一致使用相对速度的 \(\sqrt2\) 因子。

\paragraph{习题 9.12}
\(T=300.15\,\mathrm K\)，
\[
 \boxed{n=\frac p{k_BT}
 \approx3.21\times10^{17}\,\mathrm{m^{-3}}},
\quad
 \boxed{\lambda=\frac1{\sqrt2\pi d^2n}
 \approx7.79\,\mathrm m}.
\]
在摩尔质量取空气约 \(29\,\mathrm{g/mol}\)
时平均速率约 \(468\,\mathrm{m/s}\)，
因此碰撞频率
\[
 \boxed{z=\bar v/\lambda\approx60\,\mathrm{s^{-1}}}.
\]
\(\lambda\gg10^{-2}\,\mathrm m\)：
管壁碰撞将比气体分子间碰撞重要，
上式的 \(z\) 仍只是\emph{分子间}碰撞频率。

\paragraph{习题 9.13}
质量密度 \(\rho\) 对应
\(n=\rho N_A/M\)，
\[
 \boxed{n\approx3.11\times10^{16}\,\mathrm{m^{-3}}},
\quad
 \boxed{\lambda=\frac1{\sqrt2\pi d^2n}
 \approx80.3\,\mathrm m}.
\]
以 \(T=500\,\mathrm K\) 估算
\(\bar v=\sqrt{8RT/(\pi M)}
 \approx604\,\mathrm{m/s}\)，
\[
 \boxed{\tau=\lambda/\bar v\approx0.133\,\mathrm s}.
\]
高空大气真实组分及平均分子直径随高度变化，
题目只要求数量级。

\paragraph{习题 9.14}
电子远快于靶气体分子，
碰撞的相对速率近似就是电子速率，
不必引入两个\emph{同类}分子都在运动
所产生的 \(\sqrt2\) 因子。
靶分子有效碰撞截面
\(\sigma=\pi d^2\)，
单位路程碰撞次数均值 \(n\sigma\)，
所以
\[
 \boxed{\lambda_e=\frac1{n\sigma}
 =\frac1{\pi nd^2}}.
\]
若电子与靶分子都以可观速率移动，
或考虑长程库仑散射，
此简单硬球截面模型不再精确。

\paragraph{习题 9.15}
氮气在 \(T=290.15\,\mathrm K\)
的单分子平均平动、转动、总热能分别为
\[
 \boxed{\bar\epsilon_t=\frac32k_BT
 \approx6.01\times10^{-21}\,\mathrm J},\quad
 \boxed{\bar\epsilon_r=k_BT
 \approx4.01\times10^{-21}\,\mathrm J},
\]
\[
 \boxed{\bar\epsilon_{\rm thermal}=\frac52k_BT
 \approx1.00\times10^{-20}\,\mathrm J}.
\]
\(n=8.5/28\approx0.3036\,\mathrm{mol}\)，
球内氮气内能
\[
 \boxed{U=\frac52nRT\approx1.83\times10^3\,\mathrm J}.
\]
整体平移速率 \(V=65/3.6\,\mathrm{m/s}\)
对应质心\emph{轨道动能}
\[
 \boxed{K_{\rm cm}=\frac12MV^2
 \approx1.39\,\mathrm J}.
\]
气体\emph{内能}不含质心平移动能。

\paragraph{习题 9.16}
三种气体都取常温双原子理想气体模型，
均有 \(U=(5/2)nRT\)。
在 \(T=300.15\,\mathrm K\)：
\[
 \boxed{U(1\,\mathrm{mol})\approx6.24\,\mathrm{kJ}}
\]
对 \(\mathrm{H_2},\mathrm{N_2},\mathrm{O_2}\) 均相同。
各取 \(1\,\mathrm g\) 则
\(n=1/2,1/28,1/32\,\mathrm{mol}\)，
从而
\[
 \boxed{U_{\mathrm{H_2}}\approx3.12\,\mathrm{kJ}},\
 \boxed{U_{\mathrm{N_2}}\approx223\,\mathrm J},\
 \boxed{U_{\mathrm{O_2}}\approx195\,\mathrm J}.
\]
所用 \(M_{\rm H_2}=2\), \(M_{\rm N_2}=28\),
\(M_{\rm O_2}=32\,\mathrm{g/mol}\)。

\paragraph{习题 9.17}
两边初体积相同 \(V\)、压强相同 \(p\)，
初始物质的量
\(n_1=pV/(RT_1),n_2=pV/(RT_2)\)。
刚性绝热混合、无外功，两种双原子气体
定容热容均近似 \(C_{V,m}=5R/2\)，
由内能守恒
\[
 (n_1+n_2)C_{V,m}T_f
 =n_1C_{V,m}T_1+n_2C_{V,m}T_2.
\]
代入 \(n_i\propto1/T_i\)，
\[
 \boxed{T_f=\frac{2T_1T_2}{T_1+T_2}
 =\frac{2(250)(310)}{250+310}
 \approx276.8\,\mathrm K}.
\]
这并非简单算术均温 \(280\,\mathrm K\)，
因为初态相等体积、等压时两边分子数不同。

\paragraph{习题 9.18}
归一化要求
\[
 1=\int_0^{v_0}\frac{av}{v_0}\,dv
 +\int_{v_0}^{2v_0}a\,dv
 =\frac12av_0+av_0=\frac32av_0,
\]
所以
\(\boxed{a=2/(3v_0)}\)。
曲线在 \([0,v_0]\) 从零线性增至 \(a\)，
在 \([v_0,2v_0]\) 保持常数 \(a\)，
之后降为零。
\[
 \boxed{N_{v<v_0}=N/3,\qquad N_{v>v_0}=2N/3}.
\]
平均速率
\[
 \bar v=\int vf(v)dv
 =\frac{a}{v_0}\frac{v_0^3}{3}
 +a\frac{(2v_0)^2-v_0^2}{2}
 =\boxed{\frac{11}{9}v_0}.
\]
严格大于与小于的差别在连续分布
\(v=v_0\) 单点上概率为零。

\paragraph{习题 9.19}
经典方均根速率
\(v_{\rm rms}=\sqrt{3k_BT/m}\)。
取 \(m_e=9.109\times10^{-31}\,\mathrm{kg}\)，
\(m_H=1.6735\times10^{-27}\,\mathrm{kg}\)，
\(m_{\rm Na}\simeq23u\)，
得到
\[
 \boxed{v_{e,\mathrm{corona}}\approx9.54\times10^6\,\mathrm{m/s}},
\quad
 \boxed{v_{H,\mathrm{space}}\approx259\,\mathrm{m/s}},
\]
\[
 \boxed{v_{\rm Na,cold}\approx1.61\times10^{-4}\,\mathrm{m/s}}.
\]
日冕电子虽很快，仍低于光速数个百分点；
Na极低温时应注意量子气体和俘获势的适用界限。

\paragraph{习题 9.20}
逃逸速率随 \(M/R\) 的平方根变化：
\(v_{\rm esc}=v_{\rm esc,E}
 \sqrt{(M/M_E)/(R/R_E)}\)。
取地球逃逸速率约 \(11.2\,\mathrm{km/s}\)。
\[
 \boxed{v_{\rm esc,Mars}\approx5.05\,\mathrm{km/s}},
\quad
 \boxed{v_{\rm esc,Jupiter}\approx59.7\,\mathrm{km/s}}.
\]
理想气体方均根速率
\(\sqrt{3RT/M_g}\)：
\[
 \boxed{v_{\rm CO_2,Mars}\approx369\,\mathrm{m/s}},\
 \boxed{v_{\rm H_2,Mars}\approx1.73\,\mathrm{km/s}},
\]
\[
 \boxed{v_{\rm H_2,Jupiter}\approx1.27\,\mathrm{km/s}}.
\]
轻氢相对火星逃逸门槛更易受高能尾部分子逃逸；
木星引力保持轻气体更有利。
实际大气留存还依赖温度剖面、
光化学与长时日照驱动，不能仅由
\(v_{\rm rms}<v_{\rm esc}\) 推出永久稳定。
